\scanpage{160}
Soient $\underline U,\underline V$ deux \deleted{\illegible{1}} ``$\underline{\unsure{k|K}}$-algèbres'' i.e.
algèbres topologiques \deleted{\unsure{augmentées}} satisfaisant aux
conditions dites plus haut. Un hom. continu
\deleted{compatible avec l'augmentation} $\underline V\longrightarrow\underline U$
définit un \deleted{hom.}\inserted{foncteur} permis $\unsure{\mathscr C|K|\mathcal C}'^{\underline U}\longrightarrow\mathscr C'^{\underline V}$;
\deleted{et on obtient ainsi un fonct et un seul}
\deleted{\unsure{fonct}}\inserted{i.e. foncteur tel que}
\diagram{\begin{tikzcd}[ampersand replacement=\&] \mathscr C'^{\underline U}\arrow[r]\arrow[dr,"I^{\underline U}"] \& \mathscr C'^{\underline V}\arrow[d,"I^{\underline V}"] \\ {} \& \mathscr C'^{k} \end{tikzcd}}
soit commutatif.

Réciproquement, \deleted{\unsure{un}}\inserted{considérons} tel foncteur, \deleted{définit}
\deleted{\deleted{soit $v\in\underline V$.} Soit $V\in C'^{\underline U}$, on va définir une opération $\unsure{\sigma|o}_V$ dans $\underline I(W)$, \deleted{\unsure{on soit}} ceci du fait que on a une opération de $v$ ds $F(V)$ qui est un $V$-module. Les $\unsure{o|\sigma}_V$ définissent \deleted{\unsure{un}} endomorphisme du foncteur}

Comme $I^{\underline U}=I^{\underline V}\circ F$, tout endomorphisme
de $I^{\underline{\unsure{W|V}}}$ en définit un de $I^{\underline U}$, d'où
un hom de l'anneau des end.
de $I^{\underline W}$ dans celui de $I^{\underline V}$, i.e. de $\underline{\unsure{V|W}}$ dans
$\underline U$, dont on vérifie qu'il est continu.

\uline{Ainsi}, \uline{les hom continus} $\underline V\longrightarrow\underline U$
\uline{s'identifient aux hom. permis de foncteurs}.

% archive sheet number: 159
